\chapter{Decisiones arquitectónicas}
\label{sec:decisiones}

Siete ADRs documentan las decisiones que dieron forma al estado
actual. El apéndice~\ref{sec:adrs} contiene el texto completo; este
capítulo las narra.

\section{ADR-001 — Canonical API Layer}

\textbf{Problema.} Antes del rediseño, \filepath{platform/api}
dependía de \filepath{apps/cli} como librería para reusar lógica de
streaming Sentinel; y \filepath{apps/cli} dependía del core
directamente saltando la API. El resultado era un ciclo lógico
(API $\to$ CLI $\to$ core) y una bifurcación: dos caminos de
ingesta que iban a desincronizarse tarde o temprano.

\textbf{Decisión.} Extraer \code{sentinel\_streaming} a un crate
nuevo \filepath{platform/siem} del que dependan ambos; hacer que
\filepath{apps/cli} consuma \filepath{platform/api} in-proc como
cliente más.

\textbf{Consecuencia.} \filepath{platform/api} es el único lugar
con lógica; las aplicaciones son transportes. Cero ciclos.

\section{ADR-002 — FFI Isolation}

\textbf{Problema.} Un único \code{extern "C"} (bueno) pero con
RAII parcial (peligroso): \code{GmMatcher} se liberaba
manualmente. \code{GmGraph: Send + Sync} sin codificar la
invariante ``válido sólo tras finalize''. Cinco funciones C
declaradas sin uso desde Rust. Cero tests de equivalencia.

\textbf{Decisión.} Extraer el FFI a \filepath{core/matcher-ffi}.
Aplicar typestate a \code{GmGraph}. RAII completo en
\code{GmMatcher}. Añadir \code{gm_abi_version()} y chequeo en
load-time. Property tests de equivalencia DFS $\leftrightarrow$
SIMD.

\textbf{Consecuencia.} El aislamiento FFI pasó de convención a
garantía codificada. Un binario con ABI mismatch crashea al
arrancar en vez de devolver resultados sutilmente incorrectos.

\section{ADR-003 — GNN v2 Experimental}

\textbf{Problema.} El modelo \filepath{gnn_v2.onnx} se referencia
por string literal y \emph{no} está entrenado sobre todos los
datasets objetivo. Si un binario se envía con el modelo skeleton,
el scorer devuelve probabilidades sin fallar visiblemente.

\textbf{Decisión.} Aislar GNN v2 detrás de
\code{feature = "gnn-v2-experimental"}. Separar corpus export,
training e inferencia en tres componentes con contratos claros.
Añadir \emph{model gate} que rechaza modelos no validados.

\textbf{Consecuencia.} Enviar un binario a producción con GNN v2
no-validado requiere flippear explícitamente el feature flag y
pasar el gate. No es un accidente posible.

\section{ADR-004 — MCP Tools Architecture}

\textbf{Problema.} El servidor MCP original era un archivo de
\raisebox{0pt}[0pt][0pt]{\textasciitilde}1670 líneas con 54 tools
hardcoded. Cada tool llamaba al API por HTTP sin contrato tipado
entre definición de tool y endpoint, así que mismatches se
detectaban en producción.

\textbf{Decisión.} Modularizar: un archivo por tool, un registry
por dominio, una interfaz \code{Tool<I, O>} tipada con
\code{stability} y \code{version} explícitos. Tests contractuales
contra \filepath{platform/api}.

\textbf{Consecuencia.} Cambiar una tool es un PR localizado.
Deprecar una tool es una operación gobernable (marcar
\code{deprecated}, dejar existir un ciclo, luego quitar).

\section{ADR-005 — DSL Extensibility}

\textbf{Problema.} El DSL vivía en
\filepath{core/graph_hunter_core/src/dsl.rs} con entity types y
relation types fijos. Clientes pedían \code{Container},
\code{ServiceAccount}, \code{Cluster}; predicados \code{in\_subnet},
\code{matches\_regex}; operadores temporales.

\textbf{Decisión.} Extraer a \filepath{platform/dsl}. Introducir
\code{DslExtension} trait para registrar entity types, relation
types y predicados custom. Versionar gramática (\code{v1},
\code{v2-experimental}). Preservar compatibilidad de hipótesis
existentes.

\textbf{Consecuencia.} El DSL deja de ser ``hardcode en el core'' y
pasa a ser una superficie extensible. Un cliente puede pedir una
extensión sin bloquear a otros.

\section{ADR-006 — WIP Absorption}

\textbf{Problema.} Al arrancar el rediseño, \code{git status} tenía
44 archivos modificados y 14 archivos untracked: trabajo coherente
de iteraciones previas (agentic slow-lane, DLQ, invariant checker)
que no era código abandonado sino baseline no commiteado.

\textbf{Decisión.} Antes de tocar nada de la reorganización,
absorber el WIP como \emph{tres commits atómicos de baseline}. No
descartar, no posponer: integrar para poder avanzar.

\textbf{Consecuencia.} El historial del rediseño arranca de una
base coherente. Los cambios que siguen son estructurales, no
mezcla de nuevas features + mover piezas.

\section{ADR-007 — Ingestor agéntico en tres carriles}

\textbf{Problema.} La ingesta vía parsers hardcodeados en el
daemon es estática: un formato nuevo o un cambio aguas arriba
obliga a release. El parser genérico cubre casos conocidos y falla
silencioso en los demás. Ninguno aprende de su operación.

\textbf{Decisión.} Separar la ingesta en tres carriles
(cap.~\ref{sec:ingestor-agentico}): fast lane VRL compilado
determinísticamente, slow lane LLM-as-compiler offline con gate
humano, y learning lane que usa drift estructural como señal de
re-review. El contrato entre carriles es el \code{FieldConfig} +
mapping OCSF, ambos datos.

\textbf{Consecuencia.} El daemon no se recompila por formato
nuevo; se publica un nuevo mapping y el runtime lo recoge. El LLM
nunca toca datos en runtime: propone, humano aprueba, compilador
emite, fast lane ejecuta. La superficie crece
(\filepath{platform/vrl}, \filepath{platform/mapping-library},
\filepath{platform/local-llm}, tools MCP agénticas) pero la capa
core sigue intacta — shearing-layers en acción.

\section{Cómo encadenan}

Las siete ADRs no son independientes. ADR-006 es condición
necesaria para poder hacer ADR-001 (imposible reorganizar con
trabajo suelto por arriba). ADR-001 habilita ADR-002 (mover el FFI
a su propio crate requiere haber roto el ciclo). ADR-005 y ADR-004
son consecuencias de haber declarado la plataforma como superficie
versionable. ADR-003 es la aplicación del patrón \emph{model gate}
(capítulo~\ref{sec:patrones}) al caso concreto del GNN. ADR-007
reusa ese mismo patrón para el LLM del slow lane: el compilador es
opcional, remplazable, y sus fallas no afectan el fast lane.

Esto importa cuando se revisa una propuesta futura: si toca algo
cubierto por una ADR, hay que actualizar la ADR, no saltarla.
